% =====================================================================
%  Chương 1 — Mục 1 của đề: khái niệm cơ bản qua hàm và hợp hàm
%  Mã nguồn: a05/concepts.py · Notebook: notebooks/00_ham_va_hop_ham.ipynb
% =====================================================================
\chapter{Khái niệm cơ bản của CNN qua hàm và hợp hàm}
\label{ch:khainiem}

Chương này trả lời mục 1 của đề. Cách nhìn xuyên suốt là: \emph{một mạng nơ-ron tích chập
chỉ là một hàm có tham số}, và hàm đó được ghép từ nhiều hàm nhỏ,
\[
  f_\theta \;=\; f_L \circ f_{L-1} \circ \dots \circ f_2 \circ f_1 ,
  \qquad f_\theta : \mathbb{R}^{C\times H\times W} \to \mathbb{R}^{K}.
\]
Mỗi khái niệm quen thuộc của CNN (tích chập, hàm kích hoạt, pooling, softmax) là một
$f_l$ cụ thể. Huấn luyện là tìm $\theta$ làm hàm mất mát nhỏ nhất, và lan truyền ngược chỉ là
quy tắc đạo hàm của hàm hợp. Mọi hàm trong chương được viết bằng NumPy thuần trong
\texttt{a05/concepts.py}, rồi được so với PyTorch trong notebook 00. Kết quả:
\SL{concepts/checks_passed}/\SL{concepts/checks_total} phép kiểm chứng đều khớp.

\section{Neuron là một hàm affine}

Đơn vị nhỏ nhất của mạng nhận vectơ $\mathbf{x}\in\mathbb{R}^n$ và trả về một số:
\[
  z = \mathbf{w}^{\top}\mathbf{x} + b = \sum_{i=1}^{n} w_i x_i + b .
\]
Ở đây $\mathbf{w}$ và $b$ là tham số. Gom $m$ neuron thành một tầng thì được hàm affine
$\mathbf{x}\mapsto W\mathbf{x}+\mathbf{b}$ với $W\in\mathbb{R}^{m\times n}$. Mã nguồn~\ref{code:concepts__neuron}
và~\ref{code:concepts__compose} viết đúng hai ý này: một neuron, và phép hợp hàm theo thứ tự áp dụng.

\maNguon{concepts__neuron}{Neuron tuyến tính}
\maNguon{concepts__compose}{Hợp hàm theo thứ tự áp dụng}

\noindent\texttt{compose(f1, f2, f3)} trả về hàm $x\mapsto f_3(f_2(f_1(x)))$. Toàn bộ các mạng ở những
chương sau đều có dạng này, chỉ khác nhau ở danh sách hàm được truyền vào.

\subsection*{Vì sao cần hàm phi tuyến}

Ghép hai tầng affine không cho gì mới:
\[
  W_2(W_1\mathbf{x}+\mathbf{b}_1)+\mathbf{b}_2 = (W_2W_1)\,\mathbf{x} + (W_2\mathbf{b}_1+\mathbf{b}_2).
\]
Vế phải lại là một tầng affine. Nếu giữa các tầng không có hàm phi tuyến, chồng bao nhiêu tầng thì
mạng vẫn chỉ biểu diễn được hàm tuyến tính. Notebook 00 kiểm tra điều này bằng số: ghép hai tầng affine
ngẫu nhiên cho đúng kết quả của một tầng gộp, còn khi chèn ReLU vào giữa thì không còn tầng gộp nào khớp.

\maNguon{concepts__collapse_affine}{Hai tầng affine gộp thành một}

\section{Hàm kích hoạt}

Hàm kích hoạt $\sigma$ được áp dụng lên từng phần tử sau mỗi tầng affine hoặc tích chập. Ba hàm hay gặp là
\[
  \mathrm{ReLU}(z)=\max(0,z),\qquad
  \mathrm{sigmoid}(z)=\frac{1}{1+e^{-z}},\qquad
  \tanh(z)=\frac{e^{z}-e^{-z}}{e^{z}+e^{-z}} .
\]
Khi lan truyền ngược, gradient đi qua mỗi tầng bị nhân với $\sigma'(z)$. Hình~\ref{fig:kichhoat} cho thấy
đạo hàm của sigmoid không vượt quá $0{,}25$ và của tanh tiến về 0 ở hai đầu, nên qua nhiều tầng gradient bị
thu nhỏ rất nhanh. Đạo hàm của ReLU bằng đúng 1 ở mọi $z>0$, đó là lý do các CNN từ AlexNet trở đi dùng ReLU.

\begin{figure}[H]
  \centering
  \hinh{f1_kich_hoat}
  \caption{Ba hàm kích hoạt (nét liền) và đạo hàm của chúng (nét đứt)}
  \label{fig:kichhoat}
\end{figure}

\maNguon{concepts__relu}{ReLU}
\maNguon{concepts__sigmoid_grad}{Đạo hàm sigmoid viết qua chính sigmoid}

\section{Tích chập là một hàm}

\subsection{Định nghĩa}

Tầng tích chập nhận ảnh $X\in\mathbb{R}^{C\times H\times W}$ và $F$ bộ lọc $K\in\mathbb{R}^{F\times C\times k\times k}$.
Kênh ra thứ $f$ tại vị trí $(i,j)$ là
\[
  Y_{f,i,j} \;=\; b_f + \sum_{c=1}^{C}\sum_{u=0}^{k-1}\sum_{v=0}^{k-1}
      K_{f,c,u,v}\; X_{c,\; iS+u,\; jS+v},
\]
trong đó $S$ là bước trượt (stride). Công thức này thực ra là tương quan chéo (cross-correlation), không lật
bộ lọc như định nghĩa tích chập trong toán học. \texttt{nn.Conv2d} của PyTorch cũng tính đúng như vậy.
Vì bộ lọc được học từ dữ liệu, việc có lật hay không không ảnh hưởng gì.

Hình~\ref{fig:tichchap} tính tay một ví dụ của notebook 00: ảnh $5\times5$, bộ lọc $3\times3$ phát hiện
biên dọc (cột trái $+1$, cột phải $-1$).

\begin{figure}[H]
  \centering
  \hinh{f1_tich_chap}
  \caption{Tích chập tính tay: mỗi ô đầu ra là tổng tích của bộ lọc với một cửa sổ của ảnh}
  \label{fig:tichchap}
\end{figure}

Mã nguồn~\ref{code:concepts__conv2d} viết công thức trên thành hai vòng lặp theo vị trí đầu ra. Tại mỗi
vị trí, \texttt{np.tensordot} tính tổng tích giữa toàn bộ $F$ bộ lọc và cửa sổ tương ứng. Notebook 00 so
hàm này với \texttt{F.conv2d} ở ba cấu hình (stride 1 không đệm, stride 1 đệm 1, stride 2 đệm 1), và cả
ba đều khớp.

\maNguon{concepts__conv2d}{Tích chập 2D bằng NumPy}

\subsection{Padding và stride}

Đệm (padding) $P$ thêm các ô 0 quanh biên để đầu ra không bị co lại. Bước (stride) $S$ là số ô cửa sổ nhảy
mỗi lần. Kích thước đầu ra là
\[
  H_{\text{ra}} = \left\lfloor \frac{H + 2P - K}{S} \right\rfloor + 1 .
\]
Với $K=3$, chọn $P=1, S=1$ thì giữ nguyên kích thước; đây là cấu hình của mọi tầng tích chập trong
Chương~\ref{ch:mohinh}.

\begin{figure}[H]
  \centering
  \hinh{f1_pad_stride}
  \caption{Ba vị trí đầu tiên của cửa sổ $3\times3$ khi đệm $P=1$, bước $S=2$}
  \label{fig:padstride}
\end{figure}

\maNguon{concepts__out_size}{Công thức kích thước đầu ra}

\subsection{Chia sẻ trọng số: tích chập rẻ hơn tầng nối đầy đủ}

Một bộ lọc được dùng lại ở mọi vị trí của ảnh, nên số tham số của tầng tích chập không phụ thuộc vào kích
thước ảnh: nó chỉ là $C_{\text{vào}}\cdot C_{\text{ra}}\cdot k^2 + C_{\text{ra}}$. Với ảnh CIFAR $32\times32\times3$,
nối đầy đủ tới 64 đơn vị cần \SL{concepts/params/dense} tham số, còn tầng Conv $3\times3$ ra 64 kênh chỉ cần
\SL{concepts/params/conv}, ít hơn khoảng \SL{concepts/params/ratio} lần. Hai con số này được kiểm lại bằng
\texttt{numel()} của lớp PyTorch tương ứng.

\maNguon{concepts__count_params_conv}{Số tham số của tầng tích chập}

\section{Pooling}

Max pooling thay mỗi cửa sổ bằng giá trị lớn nhất của nó:
$\;Y_{c,i,j} = \max_{0\le u,v<s} X_{c,\,iS+u,\,jS+v}$. Hàm này không có tham số. Nó giảm kích thước không gian
và làm đặc trưng ít nhạy với dịch chuyển nhỏ: dù đặc trưng dịch đi một ô trong cửa sổ, giá trị lớn nhất vẫn giữ nguyên.

\begin{figure}[H]
  \centering
  \hinh{f1_pooling}
  \caption{Max pooling $2\times2$ bước 2: mỗi vùng màu cho đúng một ô ở đầu ra}
  \label{fig:pooling}
\end{figure}

\maNguon{concepts__maxpool2d}{Max pooling}

\section{Đầu ra: softmax và cross-entropy}

Sau các tầng trích đặc trưng, tensor được làm phẳng (\texttt{flatten}) rồi qua một tầng nối đầy đủ để
cho ra $K$ \emph{logits} $\mathbf{z}\in\mathbb{R}^K$. Softmax biến logits thành phân phối xác suất, và hàm mất
mát cross-entropy phạt xác suất thấp gán cho nhãn đúng $y$:
\[
  \hat{y}_k = \frac{e^{z_k}}{\sum_j e^{z_j}}, \qquad
  L(\mathbf{z}, y) = -\log \hat{y}_y = \log\sum_j e^{z_j} - z_y .
\]
Viết thẳng công thức này thì $e^{1000}$ sẽ tràn số. Trừ đi $m=\max_j z_j$ ở cả tử và mẫu không làm đổi
kết quả, nhưng đưa mọi số mũ về $\le 0$. Notebook 00 tính được
$\mathrm{softmax}(1000, 1001, 1002)$ mà không tràn, và giá trị cross-entropy khớp \texttt{F.cross\_entropy}.

\maNguon{concepts__softmax}{Softmax ổn định số}
\maNguon{concepts__cross_entropy}{Cross-entropy dùng log-sum-exp}

\section{Một CNN hoàn chỉnh là một hợp hàm}

Ghép các hàm ở trên lại, ta được một CNN nhỏ:
\[
  f = \mathrm{dense}\circ\mathrm{flatten}\circ\mathrm{maxpool}\circ\mathrm{ReLU}\circ\mathrm{conv}.
\]
Hình~\ref{fig:hopham} vẽ hợp hàm này cùng shape tensor sau từng hàm. Bảng~\ref{tab:tinytrace} là bản in thật
của notebook 00 khi chạy trên ảnh $3\times8\times8$ với bốn bộ lọc.

\begin{figure}[H]
  \centering
  \hinh{f1_hop_ham}
  \caption{CNN nhỏ viết thành hợp hàm; bốn hàm đầu trích đặc trưng, gradient đi ngược chuỗi}
  \label{fig:hopham}
\end{figure}

\begin{table}[H]
  \centering
  \caption{Shape sau từng hàm của CNN nhỏ (notebook 00)}
  \label{tab:tinytrace}
  \input{generated/tab_tiny_trace}
\end{table}

\maNguon{concepts__tiny_cnn_forward}{CNN nhỏ viết thẳng thành hợp hàm}

\noindent Hàm này nhận đúng một danh sách \texttt{(tên, hàm)} và áp dụng lần lượt, hay chính là
\texttt{compose} có kèm theo việc ghi shape. Với cùng trọng số, logits của nó khớp với
\texttt{nn.Sequential(Conv2d, ReLU, MaxPool2d, Flatten, Linear)}. Như vậy \texttt{nn.Sequential} cũng không
làm gì hơn một phép hợp hàm.

\section{Học là đạo hàm của hàm hợp}

\subsection{Quy tắc chuỗi}

Huấn luyện tìm $\theta$ làm nhỏ rủi ro thực nghiệm
$\mathcal{L}(\theta) = \frac{1}{N}\sum_{n} L(f_\theta(\mathbf{x}_n), y_n)$ bằng gradient descent
$\theta \leftarrow \theta - \eta\,\nabla_\theta\mathcal{L}$. Vì $f_\theta$ là hàm hợp, gradient theo tham số
của tầng $l$ được tính bằng quy tắc chuỗi, đi ngược từ đầu ra:
\[
  \frac{\partial L}{\partial \theta_l}
  = \frac{\partial L}{\partial \mathbf{a}_L}\,
    \frac{\partial \mathbf{a}_L}{\partial \mathbf{a}_{L-1}}\cdots
    \frac{\partial \mathbf{a}_{l+1}}{\partial \mathbf{a}_{l}}\,
    \frac{\partial \mathbf{a}_{l}}{\partial \theta_l},
  \qquad \mathbf{a}_l = f_l(\mathbf{a}_{l-1}).
\]
Mỗi tầng chỉ cần biết hai việc: nhận gradient $G = \partial L/\partial Y$ từ tầng sau, và trả về
$\partial L/\partial X$ cho tầng trước, kèm gradient theo tham số của chính nó.

\subsection{Gradient của tầng tích chập}

Vì $Y_{f,i,j}$ tuyến tính theo cả $K$ lẫn $X$, đạo hàm của nó chỉ là cộng dồn:
\[
  \frac{\partial L}{\partial K_{f,c,u,v}} = \sum_{i,j} G_{f,i,j}\, X_{c,\,iS+u,\,jS+v},\qquad
  \frac{\partial L}{\partial X_{c,p,q}} = \sum_{f}\sum_{\substack{i,j,u,v\\ iS+u=p,\; jS+v=q}} G_{f,i,j}\, K_{f,c,u,v},\qquad
  \frac{\partial L}{\partial b_f} = \sum_{i,j} G_{f,i,j}.
\]
Mã nguồn~\ref{code:concepts__conv2d_backward} duyệt lại đúng các cửa sổ mà lượt thuận đã dùng và cộng
dồn theo hai công thức đầu.

\maNguon{concepts__conv2d_backward}{Lan truyền ngược qua tầng tích chập}

\subsection{Kiểm chứng bằng sai phân hữu hạn}

Một gradient viết tay chỉ đáng tin khi đã được đối chiếu với định nghĩa đạo hàm. Hàm \texttt{grad\_check}
nhích từng phần tử $\pm\varepsilon$ và tính sai phân trung tâm $\bigl(L(\theta+\varepsilon)-L(\theta-\varepsilon)\bigr)/2\varepsilon$.
Trên ảnh $3\times9\times9$, bốn bộ lọc, stride 2, đệm 1, sai số tương đối là \SL{concepts/grad_err/K} với $\partial L/\partial K$,
\SL{concepts/grad_err/X} với $\partial L/\partial X$ và \SL{concepts/grad_err/b} với $\partial L/\partial b$, đều dưới
ngưỡng $10^{-6}$. Một bước cập nhật $K\leftarrow K-\eta\,\partial L/\partial K$ với $\eta=\SL{concepts/gd_step/eta}$ làm
$L$ giảm từ \SL{concepts/gd_step/before} xuống \SL{concepts/gd_step/after}, đúng hướng mà gradient chỉ ra.

\maNguon{concepts__grad_check}{Kiểm tra gradient bằng sai phân trung tâm}

\begin{phathien}{Tóm tắt Chương 1}
Mỗi thành phần của CNN là một hàm: tích chập là hàm tuyến tính có chia sẻ trọng số, ReLU là hàm phi tuyến
đặt giữa các tầng tuyến tính để chúng không gộp lại thành một tầng, pooling là hàm lấy cực đại cục bộ,
softmax và cross-entropy biến logits thành mất mát. CNN là hợp của các hàm đó, và việc học là tính đạo hàm
của hàm hợp bằng quy tắc chuỗi. Toàn bộ \SL{concepts/checks_total} hàm viết tay khớp với PyTorch.
\end{phathien}
